% RESEARCH_DRAFT_SHA256: 5353efca5eb87e6b3d964668a9676375dbfcc74e2670e5ca2e56c1e8738fe43c
% Layout and source-reference wrappers expanded for the web renderer.
\documentclass[11pt]{article}
\usepackage[margin=25mm]{geometry}
\usepackage{amsthm}
\usepackage{fontspec,unicode-math}
\setmainfont{Latin Modern Roman}
\setmathfont{Latin Modern Math}
\usepackage{amsmath,amssymb,mathtools,xurl,hyperref}
\hypersetup{colorlinks=true,urlcolor=blue,linkcolor=blue}
\setcounter{secnumdepth}{0}
\setlength{\parindent}{0pt}
\setlength{\parskip}{5pt}
\setlength{\emergencystretch}{3em}
\newcommand{\sref}[2]{\hyperlink{source:#1:#2}{[S#2]}}
\newcommand{\cataloguescope}[1]{\paragraph{Source record.}#1}
\newcommand{\floor}[1]{\left\lfloor#1\right\rfloor}
\newtheorem{proposition}{Proposition}
\newtheorem{lemma}[proposition]{Lemma}
\begin{document}
\paragraph{Research manuscript.} OpenAI \texttt{gpt-6-astra}, reasoning \texttt{ultra}; 11 October 2026.
\paragraph{Preserved source body.} The complete catalogue mathematical text and sources follow with layout and source links expanded, including source qualifications and its historical preparation notes. New research begins at the section headed Research attempt.

% UnsolvedMath ID 30005831; source catalogue code OWR-14298166-008
\section{Lower Bounds for Polynomial Calculus with Inverses}\label{30005831}
{\small UnsolvedMath ID 30005831 · Computer Science}
\subsection{Definitions and mathematical statement}

\paragraph{Shared notation.}Standard number systems are $\mathbb N,\mathbb Z,\mathbb Q,\mathbb R,\mathbb C$. Specialized notation and index ranges follow the cited source and the definitions below.


Polynomial time means a number of elementary steps bounded by a polynomial in the encoded input length. A decision problem is in $\mathrm{NP}$ if yes-instances have polynomial-length certificates verifiable in polynomial time; NP-hardness is under polynomial-time reductions.

A refutation is a derivation of contradiction from an unsatisfiable constraint system using specified inference rules. Proof size counts the encoded derivation, and degree in algebraic proof systems is the largest polynomial degree appearing.
\paragraph{Record scope and qualifications.}
The catalogue title is retained as metadata. Named certificates, contracts, bounds or frameworks in that title are not asserted as results of the original source extract.
Source-record limitation: The exact inversion rule, underlying field, permitted nonzero assumptions and representation of tensor-isomorphism equations are missing; ordinary Polynomial Calculus definitions cannot supply the stronger PC+Inv target.
\paragraph{Statement recorded in the supplied source.}
Can polynomial calculus with inversion solve Tensor Isomorphism in polynomial time? More generally, prove lower bounds for polynomial calculus with inversion on any system of equations. \href{https://doi.org/10.4171/owr/2024/15}{[S2]}
\paragraph{Source status.}
The supplied catalogue labels this source record open; that metadata was checked against the supplied snapshot on 2026-10-11. The label is preserved for provenance and is not an independent verification of present open status or mathematical correctness. The catalogue literature assessment is dated 2026-08-22.
\subsection{Short English statement}
The source asks about lower Bounds for Polynomial Calculus with Inverses.
\subsection{Sources}
\begin{itemize}
\item[S1] Ulam AI and the UnsolvedMath Contributors, UnsolvedMath: A Curated Collection of Open Mathematics Problems, record 30005831 (OWR-14298166-008). Catalogue curation and metadata: CC BY 4.0. The source extract is retained; definitions, English summaries and qualifications are editorial additions.
\url{https://huggingface.co/datasets/ulamai/UnsolvedMath}.
\item[S2] Proof Complexity and Beyond (2024), Oberwolfach Reports, DOI 10.4171/owr/2024/15, pp. 871--934
\url{https://doi.org/10.4171/owr/2024/15}.
\item[S3] Galesi et al., On the Algebraic Proof Complexity of Tensor Isomorphism
\url{https://doi.org/10.4230/LIPIcs.CCC.2023.4}.
\item[S4] Oberwolfach Report 15/2024
\url{https://oa.tib.eu/renate/bitstreams/b0ef46dc-1bdf-45b7-a4f5-ca0a9e60a7d6/download}.
\item Source-specific Oberwolfach Reports licences, including CC BY-SA 4.0 where applicable, are retained; see THIRD\_PARTY\_NOTICES.md.
\end{itemize}
\paragraph{Preparation and review.}
The original source and UnsolvedMath contributors retain attribution for the source extract. Codex (GPT-6; exact serving revision not exposed) assisted with definitions, English summaries, scope notes and formatting. No new human mathematical authorship or independently verified solution is asserted. Permanent source ID 30005831 and code OWR-14298166-008 checked against the complete supplied registry and current repository. Distinct upstream records, including repeated questions and broad or damaged extracts, are retained for complete coverage of the source CS category. A repeated record is not a new independent conjecture. Complete mathematical-equivalence and current-literature review has not been performed. Preparation and final review follow contribution prompts C and D, with source gaps disclosed. Structural and rendering checks do not establish a new theorem or open status.
\label{end:30005831}

\section{Research attempt}
\subsection{Literature, source context and scope}
The full original catalogue body and source qualifications are preserved above. The relevant primary workshop passages were read in the cached publisher report. Historical claims and subsequent attributed results are distinguished from the mathematical work below. A restricted example, standard lemma or known theorem is not counted as new progress on the entire original target. Missing assumptions and repeated or rhetorical extracts are not silently turned into different solved problems.

\subsection{Recovering the actual inversion rule}
The primary paper defines the additional rule for square polynomial matrices: $AB=I$ permits $BA=I$. It is not merely adjoining scalar reciprocals. Its Lemma3.1 already supplies the following rank-principle route; we rework it as a failed candidate for a lower bound, not new progress.
\begin{proposition}
For $1\le r<n$, the equations $XY=I_n$, with $X$ of size $n\times r$ and $Y$ of size $r\times n$, have a degree-three refutation in $\mathrm{PC}+\mathrm{Inv}$.
\end{proposition}
\begin{proof}
Partition $X=(X_0;X_1)$ and $Y=(Y_0\ Y_1)$ with square $r\times r$ blocks $X_0,Y_0$. The premises give $X_0Y_0=I_r$, $X_0Y_1=0$, and $X_1Y_1=I_{n-r}$. Inversion yields $Y_0X_0=I_r$. Entrywise variable multiplication and linear combination give
\[
Y_0(X_0Y_1)-(Y_0X_0-I_r)Y_1=Y_1=0
\]
in degree at most three. Multiplying these linear equations by entries of $X_1$ yields $X_1Y_1=0$. Subtract a diagonal equation of $X_1Y_1=I_{n-r}$ to obtain $1=0$. Every matrix multiplication here expands into polynomially many permitted variable multiplications and sums.
\end{proof}
Thus an ordinary-PC rank lower bound does not transfer after adding this rule.
\subsection{A scope check that prevents an unsound rule extension}
For square matrices over a field, soundness follows directly from finite-dimensional linear algebra: $AB=I$ makes $A$ surjective, so it has zero kernel and is bijective; multiplying by its inverse gives $B=A^{-1}$ and $BA=I$. This also checks polynomial substitution instances pointwise at any assignment over an extension field. Square shape is essential. For $A=(1\ 0)$ and $B=(1\ 0)^T$, one has $AB=I_1$ but $BA=\operatorname{diag}(1,0)\ne I_2$. Treating the original rectangular rank equations themselves as an inversion-rule instance would be invalid; the proof above applies the rule only to their square corner.
\subsection{Verification and remaining gaps}
The degree-three identity was expanded in independent commuting variables for $n=3,r=2$, with exact cancellation of all cubic terms, and the rectangular counterexample was checked exactly. The proof is an attributed known auxiliary upper bound. It does not bound arbitrary inversion substitutions, produce a hard tensor-isomorphism family for the strengthened system, or establish a general lower bound. Those are the unresolved original targets; an interpretation using a different scalar-inversion system would answer a different question.

\section{References and source attribution}
Primary workshop report: \url{https://doi.org/10.4171/owr/2024/15}. Relevant passages read: 2024/15 p.909, PC+Inv and tensor-isomorphism lower-bound questions.

\url{https://drops.dagstuhl.de/storage/00lipics/lipics-vol264-ccc2023/LIPIcs.CCC.2023.4/LIPIcs.CCC.2023.4.pdf}. Read primary pp.4--5 definition/conjecture and p.10 Lemma3.1 rank-principle proof; the reworked auxiliary derivation is explicitly attributed.

All retrievals were checked on 11 October 2026. Original question attribution remains in the preserved Sources section. External theorems are attributed and not claimed as new discoveries or fully verified merely by citation.


\paragraph{Integration review.} Codex (GPT-6; exact serving revision not exposed) reviewed the written arguments and expanded layout and source-reference wrappers. Repository metadata and the production web renderer were checked using Python 3.13 and JavaScript with MathJax 3.2.2. This model-assisted review is not independent human proof verification; no PDF compilation or formal proof checking is asserted.

\section{Completion Estimate}
\noindent\textbf{COMPLETION: 0\%}
No substantiated new progress toward the entire original target is claimed. The worked results and failed route are delimited above; a known result, illustrative model, or auxiliary obstruction is not counted as original-target completion.

\section{Final status and attribution}
\textbf{Attempt status: unresolved.} This describes this research attempt, not a new ruling on the source catalogue's status or on external literature claims.
The mathematics and this self-audit were generated by OpenAI \texttt{gpt-6-astra}, reasoning \texttt{ultra}, on 11 October 2026. Codex (GPT-6; exact serving revision not exposed) orchestrated the batch, read and reviewed its written mathematical arguments, and contributed formatting/repository review. No human mathematical authorship, independent human proof checking, or formal verification is claimed. The historical preparation paragraphs in the reproduced catalogue body concern that body only.

Contribution prompt A was applied to the complete definition and supplied source record, followed by the adversarial checks of prompt B. Local reproduction files are \texttt{batch025/checks.py}, \texttt{computations.json}, and the per-ID audit JSON. They were not public links at preparation time. Numerical checks are finite and do not prove unrestricted results. No LaTeX engine was found locally; compilation was not performed by this research agent.

\end{document}
